%==============================================================================
% macros/notation.tex — 统一项目符号表
%
% 设计原则：
%   1. 与 Docs/black/path_A_method_skeleton.md §3.1 Notation 表保持一致
%   2. 加宏的目的：(a) 一处改名全文同步 (b) 防止 reviewer 抱怨符号漂移
%   3. 所有专业符号都通过宏，不直接写裸 LaTeX
%
% 修改时务必同步更新 path_A_method_skeleton.md
%==============================================================================

% ----- 数集 -----
\newcommand{\R}{\mathbb{R}}
\newcommand{\N}{\mathbb{N}}
\newcommand{\Z}{\mathbb{Z}}
\newcommand{\E}{\mathbb{E}}
\newcommand{\Prob}{\mathbb{P}}

% ----- 概率 / 测度 -----
\newcommand{\Law}{\mathrm{Law}}
\newcommand{\KL}{\mathrm{KL}}
\newcommand{\Wass}[1]{W_{#1}}                 % e.g. \Wass{1}, \Wass{2}
\newcommand{\TV}{\mathrm{TV}}
\newcommand{\BV}{\mathrm{BV}}
\newcommand{\Ent}{\mathrm{Ent}}
\newcommand{\Lip}{\mathrm{Lip}}

% ----- 分布与路径 -----
\newcommand{\rhotrue}{\rho^{\star}}             % 目标 PDE 解分布
\newcommand{\rhotau}{\rho_{\tau}}               % 扩散时间 tau 处的分布
\newcommand{\ptau}{p_{\tau}}                    % noised density
\newcommand{\score}{s}                          % score
\newcommand{\sth}{s_{\theta}}                   % score network
\newcommand{\Dth}{D_{\theta}}                   % denoiser network
\newcommand{\sigtau}{\sigma(\tau)}              % 噪声尺度
\newcommand{\Gtau}{G_{\tau}}                    % heat kernel

% ----- PDE -----
\newcommand{\dt}{\partial_{t}}
\newcommand{\dx}{\partial_{x}}
\newcommand{\dtau}{\partial_{\tau}}
\newcommand{\Lap}{\Delta}
\newcommand{\flux}{f}                           % 通量函数 f(u)
\newcommand{\physvis}{\nu_{\mathrm{phys}}}      % PDE 物理黏性
\newcommand{\snr}{\mathrm{SNR}}
\newcommand{\amp}{\Lambda}                      % Score Shocks 的指数放大因子

% ----- 损失 -----
\newcommand{\Ldsm}{\mathcal{L}_{\mathrm{DSM}}}
\newcommand{\Lpde}{\mathcal{L}_{\mathrm{PDE}}}
\newcommand{\Lent}{\mathcal{L}_{\mathrm{ent}}}
\newcommand{\Lbv}{\mathcal{L}_{\mathrm{BV}}}
\newcommand{\Lburg}{\mathcal{L}_{\mathrm{Burg}}}
\newcommand{\Rent}{\mathrm{R}_{\mathrm{ent}}}

% ----- 集合 / 子集 -----
\newcommand{\Shockphys}{\Sigma_{\mathrm{phys}}}
\newcommand{\Shockscore}{\Sigma_{\mathrm{score}}}

% ----- 物理空间解与轨迹（SYMBOL.md §5/§10） -----
\newcommand{\physsol}{\mathbf{u}^{\star}}        % Kruzhkov 熵解
\newcommand{\physsolvis}{\mathbf{u}^{\nu}}       % 粘性 PDE 解
\newcommand{\initdat}{\mathbf{u}_{0}}            % PDE 初值
\newcommand{\rhotphys}{\rho_{t}}                 % 物理时间 t 处分布

% ----- Score 残差与误差 -----
\newcommand{\errsc}{e}                           % score 残差 e = s_theta - s

% ----- 几何对象扩展 -----
\newcommand{\dShock}{d_{\Sigma}}                 % 到 shock 流形的符号距离

% ----- BV-aware 参数化分量（SYMBOL.md §9） -----
\newcommand{\phisbg}{\phi^{\mathrm{sm}}_{\theta}}  % 平滑背景势函数
\newcommand{\phissh}{\phi^{\mathrm{sh}}_{\theta}}  % shock 符号距离场
\newcommand{\jumpamp}{\kappa_{\theta}}           % 局部跳跃强度

% ----- 轨迹与生成样本（SYMBOL.md §10） -----
\newcommand{\unat}{u^{\natural}}                 % 真实 score 驱动的反向 ODE 轨迹
\newcommand{\uhat}{\widehat{u}}                  % 网络 score 驱动的反向 ODE 轨迹
\newcommand{\uth}{u^{\theta}}                    % 最终生成样本 = \uhat(0)
\newcommand{\muth}{\mu_{\theta}}                 % 生成分布 = Law(u^theta)

% ----- 实用工具 -----
\newcommand{\abs}[1]{\lvert #1 \rvert}
\newcommand{\norm}[1]{\lVert #1 \rVert}

% ----- 排版小工具 -----
\newcommand{\eg}{e.g.,\ }
\newcommand{\ie}{i.e.,\ }
\newcommand{\cf}{cf.\ }

% ----- 备注高亮（写作期临时使用，提交前删除）-----
\newcommand{\todo}[1]{\textcolor{red}{\textbf{TODO: }#1}}
\newcommand{\note}[1]{\textcolor{blue}{[Note: #1]}}
